% SHARD-ID: AQM-39-TWO-QUTRIT-STABILIZER-DESCENT-COUNTEREXAMPLE
% SHARD-TITLE: Two-Qutrit Stabilizer Descent Counterexample
% SHARD-SUMMARY: Gives an explicit global Lagrangian stabilizer with nonzero descent defect.
% SHARD-SUMMARY: Shows that the associated two-qutrit stabilizer state has trivial one-site stabilizer restrictions.
% SHARD-SUMMARY: Records the corrected local-to-global statement for Lagrangian stabilizers.
% SHARD-KEYWORDS: qutrit, stabilizer, Lagrangian, descent defect, entanglement, local-to-global

\section{Two-Qutrit Stabilizer Descent Counterexample}
\label{sec:two-qutrit-stabilizer-descent-counterexample}

\subsection{Status and Source Anchors}

This shard is a checked finite calculation using the convention and theorems
of AQM-38.  It is the promised counterexample to the overly strong claim that
choosing a global Lagrangian stabilizer is the same thing as satisfying a
local-to-global principle.  The Weyl commutation convention is the same
prime-field convention as AQM-08, AQM-22, and \texttt{CONVENTIONS.md} (l).

\subsection{Lamport Dependency Skeleton}

\[
\begin{array}{rcl}
D1&:&U=A\sqcup B,\quad E(U)=\mathbb F_3^2\oplus\mathbb F_3^2.\\
D2&:&\omega(e,e')=z_Ax_A'-x_Az_A'+z_Bx_B'-x_Bz_B'.\\
D3&:&\ell_1=(1,0,-1,0),\quad \ell_2=(0,1,0,1),\quad
      L=\langle\ell_1,\ell_2\rangle.\\
L1&:&L\text{ is Lagrangian}.\\
L2&:&L\cap E(A)=0\text{ and }L\cap E(B)=0.\\
T1&:&D_U(L;\{A,B\})\cong L\ne0.\\
T2&:&\text{The stabilizer state is entangled and has maximally mixed
      one-site marginals}.\\
T3&:&\text{A global Lagrangian is not the same as stabilizer descent.}
\end{array}
\]

\subsection{The Label Calculation D1--T1}

Let \(U=A\sqcup B\), with one qutrit on each site.  Write labels as
\[
  (x_A,z_A,x_B,z_B)\in
  \mathbb F_3^2\oplus\mathbb F_3^2
\]
and use the direct-sum symplectic form
\[
  \omega(e,e')=
  z_Ax_A'-x_Az_A' + z_Bx_B'-x_Bz_B' .
\]
Define
\[
  \ell_1=(1,0,-1,0),\qquad
  \ell_2=(0,1,0,1),
  \qquad
  L=\langle \ell_1,\ell_2\rangle .
\]

\paragraph{Lemma \(L1\).}
The subspace \(L\) is Lagrangian.

\emph{Proof.}
The only nontrivial generator pairing is
\[
  \omega(\ell_1,\ell_2)
  =
  0\cdot0-1\cdot1+0\cdot0-(-1)\cdot1
  =
  -1+1=0 .
\]
Same-generator pairings vanish because \(\omega\) is alternating, so \(L\) is
isotropic.  The ambient symplectic vector space has dimension \(4\), and
\(\ell_1,\ell_2\) are independent, so \(\dim L=2\).  Since the form is
nondegenerate, \(\dim L+\dim L^\perp=4\), so \(\dim L^\perp=2\).  Isotropy
gives \(L\subseteq L^\perp\), and equal dimensions give \(L=L^\perp\).  Hence
\(L\) is Lagrangian.

\paragraph{Lemma \(L2\).}
\[
  L\cap E(A)=0,\qquad L\cap E(B)=0.
\]

\emph{Proof.}
A general element of \(L\) is
\[
  a\ell_1+b\ell_2=(a,b,-a,b).
\]
If it is supported only on \(A\), then \(-a=0\) and \(b=0\), hence
\(a=b=0\).  If it is supported only on \(B\), then \(a=0\) and \(b=0\).  These
are exactly the displayed intersections.

\paragraph{Theorem \(T1\).}
For the disjoint cover \(\{A,B\}\),
\[
  D_U(L;\{A,B\})\cong L.
\]
In particular the defect has \(\mathbb F_3\)-dimension \(2\).

\emph{Proof.}
By Lemma \(L2\),
\[
  \operatorname{Loc}_U(L;\{A,B\})
  =
  (L\cap E(A))+(L\cap E(B))=0.
\]
The definition from AQM-38 gives
\[
  D_U(L;\{A,B\})
  =
  L/\operatorname{Loc}_U(L;\{A,B\})
  =
  L.
\]

\subsection{The Stabilizer State T2}

\paragraph{Theorem \(T2\).}
With the standard qutrit Weyl lift, the stabilizer labels \(L\) define a full
global stabilizer state, but no nonzero local stabilizer label on either
single site is inherited from \(L\).

\emph{Proof.}
The two generator operators are
\[
  W(\ell_1)=X_AX_B^{-1},\qquad W(\ell_2)=Z_AZ_B .
\]
The Weyl multiplier is trivial on all of \(L\), because
\[
  (b,b)\cdot(a',-a')=0
\]
for \(a\ell_1+b\ell_2\) and \(a'\ell_1+b'\ell_2\).  Thus the phase
\(\alpha=1\) satisfies the AQM-38 phase condition.  The two displayed
operators commute by Lemma \(L1\).  The vector
\[
  |\Phi_-\rangle=\frac1{\sqrt3}\sum_{s\in\mathbb F_3}|s,-s\rangle
\]
is fixed by both operators.  Indeed, \(X_AX_B^{-1}\) sends
\(|s,-s\rangle\) to \(|s+1,-s-1\rangle=|s+1,-(s+1)\rangle\), only
reindexing the sum.  Also \(Z_AZ_B\) multiplies \(|s,-s\rangle\) by
\(\chi(s-s)=1\).  Since \(L\) is Lagrangian, AQM-38 Theorem \(T1\) says the
common stabilizer space is one-dimensional, so this is the full stabilizer
state.

The one-site density matrix on \(A\) is
\[
  \operatorname{Tr}_B(|\Phi_-\rangle\langle\Phi_-|)
  =
  \frac13\sum_{s\in\mathbb F_3}|s\rangle\langle s|
  =
  I_3/3,
\]
and the same computation gives \(I_3/3\) on \(B\).  Thus the one-site
restrictions are maximally mixed.  This agrees with Lemma \(L2\): the global
Lagrangian has no nonzero labels supported on a single site.

\subsection{Corrected Local-to-Global Statement T3}

\paragraph{Theorem \(T3\).}
Choosing a full stabilizer is not equivalent to local-to-global gluing.  The
correct label-level statement is:
\[
\begin{array}{c}
\text{full global stabilizer state}\\
\Longleftrightarrow\\
\text{global phased Lagrangian }L\leq E(U),
\end{array}
\]
whereas
\[
\begin{array}{c}
\text{full global stabilizer generated from the cover }\mathcal U\\
\Longleftrightarrow\\
\text{global phased Lagrangian }L\leq E(U)\text{ with }
D_U(L;\mathcal U)=0,
\end{array}
\]
plus phase compatibility on overlaps.

\emph{Proof.}
The first equivalence is AQM-38 Theorem \(T1\).  The second equivalence is
AQM-38 Theorem \(T2\) applied to a Lagrangian \(L\), with phase compatibility
added because labels alone do not choose eigenvalues.  AQM-38 Theorem \(T3\)
proves that product local stabilizer states satisfy the condition.  Theorem
\(T1\) above proves that the Lagrangian condition alone does not imply it.

\subsection{Conclusion}

The broad equivalence must be rejected in this form.  A Lagrangian is a full
global stabilizer.  A zero descent defect is the extra condition that this
stabilizer is generated from the local pieces of a chosen cover.  Nonzero
defect is the stabilizer-label analogue of the state-presheaf phenomenon from
AQM-27: local restrictions can miss global correlations.
